% GENERATED FILE. Edit canonical lessons, then run npm run generate:books.
% canonical-source-hash: fnv1a64:9e9bcb5ee742a955
% canonical-lessons: TE-C319-vellali, TE-C319-otu, TE-S175-letter-oo, TE-C319-nadavadam-valla, TE-R319-advice-and-reasons

\chapter{Must, Must Not, and Because}
\label{ch:te-must-must-not-and-because}
\hlchaptermodality{319}

\begin{chapteropening}
\textbf{By the end of this chapter:} \emph{I can say what someone must, must not or had better do with -āli, -akūḍadu and -tē mañcidi, and explain why with -aḍaṁ valla, endukaṇṭē and kābaṭṭi.}

Complete TE-R319-advice-and-reasons: give advice and duties, put a reason before or after its result, and read the new letter \te{ఓ}.
\end{chapteropening}

\section[\texorpdfstring{\te{వెళ్ళాలి} (veḷḷāli)}{వెళ్ళాలి (veḷḷāli)}]{\te{వెళ్ళాలి} (veḷḷāli) — must, must not, had better}

\label{lesson:TE-C319-vellali}



\begin{quote}
Say \textquotedblleft{}I want a ticket\textquotedblright{}. (\emph{nāku oka ṭikkeṭṭu kāvāli}.) Then \textquotedblleft{}It is raining, so I am at home.\textquotedblright{} (\emph{varṣaṁ paḍutōndi kābaṭṭi nēnu iṇṭlō unnānu}.)
\end{quote}

\subsection*{You'll want to know: -\te{ఆలి}, must}
\begin{quote}
\textbf{\te{నేను} \te{ఇప్పుడు} \te{వెళ్ళాలి}.}
\end{quote}

\begin{quote}
\emph{nēnu ippuḍu veḷḷāli.} — \textquotedblleft{}I must go now.\textquotedblright{}
\end{quote}

\begin{quote}
\textbf{\te{మీరు} \te{ఈ} \te{మందు} \te{ప్రతి} \te{రోజు} \te{తీసుకోవాలి}.}
\end{quote}

\begin{quote}
\emph{mīru ī mandu prati rōju tīsukōvāli.} — \textquotedblleft{}You should take this medicine every day.\textquotedblright{}
\end{quote}

\textbf{-\te{ఆలి}} sits on the verb's base: \textbf{\te{వెళ్ళు}} $\to$ \textbf{\te{వెళ్ళాలి}}, \textbf{\te{చేయు}} $\to$ \textbf{\te{చేయాలి}}, \textbf{\te{తిను}} $\to$ \textbf{\te{తినాలి}}, \textbf{\te{చదువు}} $\to$ \textbf{\te{చదవాలి}}, and \textbf{\te{వచ్చు}} $\to$ \textbf{\te{రావాలి}}. It does not change for the person: \textbf{\te{నేను}}, \textbf{\te{నువ్వు}}, \textbf{\te{అతను}}, all \textbf{\te{వెళ్ళాలి}}. It covers both \textquotedblleft{}must\textquotedblright{} and \textquotedblleft{}should\textquotedblright{}.

You have said this ending for a long time. \textbf{\te{కావాలి}} is \textbf{-\te{ఆలి}} on \textbf{\te{కా}-}, \textquotedblleft{}to become\textquotedblright{}: \textquotedblleft{}it must come to be (mine)\textquotedblright{}, so \textquotedblleft{}it is needed, it is wanted\textquotedblright{}.

\subsection*{You'll want to know: -\te{అకూడదు}, must not}
\begin{quote}
\textbf{\te{జ్వరం} \te{ఉన్నప్పుడు} \te{బడికి} \te{వెళ్ళకూడదు}.}
\end{quote}

\begin{quote}
\emph{jvaraṁ unnappuḍu baḍiki veḷḷakūḍadu.} — \textquotedblleft{}When you have a fever, you must not go to school.\textquotedblright{}
\end{quote}

\begin{quote}
\textbf{\te{ఇప్పుడు} \te{కాఫీ} \te{తాగకూడదు}.}
\end{quote}

\begin{quote}
\emph{ippuḍu kāphī tāgakūḍadu.} — \textquotedblleft{}You must not drink coffee now.\textquotedblright{}
\end{quote}

\textbf{-\te{అకూడదు}}: \textbf{\te{వెళ్ళకూడదు}}, \textbf{\te{తాగకూడదు}}, \textbf{\te{చేయకూడదు}}. \textbf{\te{ఉన్నప్పుడు}} (\emph{unnappuḍu}, \textquotedblleft{}when there is\textquotedblright{}) is \textbf{\te{ఉన్న}} + \textbf{\te{అప్పుడు}}, built like \textbf{-\te{తున్నప్పుడు}}.

\subsection*{You'll want to know: -\te{తే} \te{మంచిది}, had better}
\begin{quote}
\textbf{\te{మీరు} \te{ఆసుపత్రికి} \te{వెళ్తే} \te{మంచిది}.}
\end{quote}

\begin{quote}
\emph{mīru āsupatriki veḷtē mañcidi.} — \textquotedblleft{}You had better go to the hospital.\textquotedblright{}
\end{quote}

Word for word, \textquotedblleft{}if you go, it is good\textquotedblright{}: advice, softer than \textbf{-\te{ఆలి}}. The \textbf{-\te{తే}}, \textquotedblleft{}if\textquotedblright{}, gets a lesson of its own later in this core; here keep \textbf{-\te{తే} \te{మంచిది}} as one piece.

\subsection*{Your turn}
Say these aloud:

\begin{itemize}
\raggedright
  \item \textquotedblleft{}I must go now\textquotedblright{} — \emph{nēnu ippuḍu veḷḷāli}
  \item \textquotedblleft{}You should take this medicine every day\textquotedblright{} — \emph{mīru ī mandu prati rōju tīsukōvāli}
  \item \textquotedblleft{}You must not drink coffee now\textquotedblright{} — \emph{ippuḍu kāphī tāgakūḍadu}
  \item \textquotedblleft{}You had better go to the hospital\textquotedblright{} — \emph{mīru āsupatriki veḷtē mañcidi}
\end{itemize}

\begin{tcolorbox}[breakable,colback=teal!4,colframe=teal!35!black,title={Before you move on}]
\textquotedblleft{}Must come\textquotedblright{}? (\textbf{\te{రావాలి}}.) \textquotedblleft{}Must not eat\textquotedblright{}? (\textbf{\te{తినకూడదు}}.) \textquotedblleft{}Had better go\textquotedblright{}? (\textbf{\te{వెళ్తే} \te{మంచిది}}.) Which word you knew already ends in \textbf{-\te{ఆలి}}? (\textbf{\te{కావాలి}}.)
\end{tcolorbox}
\section[\texorpdfstring{\te{ఓటు} (ōṭu)}{ఓటు (ōṭu)}]{\te{ఓటు} (ōṭu) — a vote}

\label{lesson:TE-C319-otu}



\begin{quote}
Before the new one: say \textquotedblleft{}I must go\textquotedblright{} and \textquotedblleft{}you must not drink coffee now\textquotedblright{}. (\emph{nēnu veḷḷāli}; \emph{ippuḍu kāphī tāgakūḍadu}.)
\end{quote}

\subsection*{You'll want to know: \te{ఓటు}}
\textbf{\te{ఓటు}} — \emph{ōṭu} — \textquotedblleft{}a vote\textquotedblright{}.

\begin{quote}
\textbf{\te{ప్రతి} \te{ఒక్కరూ} \te{ఓటు} \te{వేయాలి}.}
\end{quote}

\begin{quote}
\emph{prati okkarū ōṭu vēyāli.} — \textquotedblleft{}Everyone must vote.\textquotedblright{}
\end{quote}

\begin{quote}
\textbf{\te{నేను} \te{నిన్న} \te{ఓటు} \te{వేశాను}.}
\end{quote}

\begin{quote}
\emph{nēnu ninna ōṭu vēśānu.} — \textquotedblleft{}I voted yesterday.\textquotedblright{}
\end{quote}

A vote is \textquotedblleft{}put\textquotedblright{}: \textbf{\te{ఓటు} \te{వేయు}} (\emph{ōṭu vēyu}, \textquotedblleft{}to cast a vote\textquotedblright{}), with \textbf{\te{వేయు}}, \textquotedblleft{}to put\textquotedblright{}. \textbf{\te{ప్రతి} \te{ఒక్కరూ}} (\emph{prati okkarū}, \textquotedblleft{}every single person\textquotedblright{}) is \textbf{\te{ప్రతి}}, \textquotedblleft{}every\textquotedblright{}, with a word for \textquotedblleft{}one person\textquotedblright{}.

Its first letter, a long \emph{ō} at the start of a word, is one you have not written. The next lesson takes it out of this word.

\begin{cousinweb}
\textbf{\te{ఓటు}} (\emph{ōṭu}) is a loan from English \emph{vote}, with the \textbf{-\te{ఉ}} that Telugu gives most borrowed nouns, as in \textbf{\te{బస్సు}} and \textbf{\te{టిక్కెట్టు}}. The English \emph{v} drops away, the vowel stays long, and the English \emph{t} becomes \textbf{\te{ట}}, as it does in \textbf{\te{టిక్కెట్టు}}.
\end{cousinweb}

\subsection*{Your turn}
Say these aloud:

\begin{itemize}
\raggedright
  \item \emph{ōṭu}, with a long \emph{ō}
  \item \textquotedblleft{}Everyone must vote\textquotedblright{} — \emph{prati okkarū ōṭu vēyāli}
  \item \textquotedblleft{}I voted yesterday\textquotedblright{} — \emph{nēnu ninna ōṭu vēśānu}
\end{itemize}

\begin{tcolorbox}[breakable,colback=teal!4,colframe=teal!35!black,title={Before you move on}]
What does \textbf{\te{ఓటు}} mean? (\textbf{A vote.}) \textquotedblleft{}To cast a vote\textquotedblright{}? (\textbf{\te{ఓటు} \te{వేయు}}.) Where does the word come from? (\textbf{English \emph{vote}.})
\end{tcolorbox}
\section[\texorpdfstring{\te{ఓ} (ō)}{ఓ (ō)}]{\te{ఓ} (ō) — the letter that opens \te{ఓటు}}

\label{lesson:TE-S175-letter-oo}



\begin{quote}
Before the new one: write \textbf{\te{ఒ}} and \textbf{\te{ో}} once each, and name them. (Short \emph{o} at the start of a word; long \emph{ō} after a consonant, as in \textbf{\te{తోలు}}.)

Now say \textbf{\te{ఓటు}} (\emph{ōṭu}), \textquotedblleft{}a vote\textquotedblright{}. This lesson takes its first letter out and puts it in your hand.
\end{quote}

\subsection*{Script you'll notice: \te{ఓ}}
\textbf{\te{ఓ}} — \emph{ō}, a long \emph{ō} at the start of a word.

After a consonant, long \emph{ō} is the sign \textbf{\te{ో}}. At the start of a word it stands as a letter of its own, \textbf{\te{ఓ}}, and it pairs with the short \textbf{\te{ఒ}} you know, the way \textbf{\te{ఏ}} pairs with \textbf{\te{ఎ}}.

In print \textbf{\te{ఓ}} is \textbf{\te{ఒ}} with a small hooked stroke added at its top, rising and curling over to the right. Say the pair aloud: \textbf{\te{ఒ}} \emph{o}, short; \textbf{\te{ఓ}} \emph{ō}, held longer.

\subsection*{Writing: \te{ఓ} — copy the printed shape}
Put your pen beside \textbf{\te{ఓ}} and copy it, slowly, and larger than it is printed. Write the body of \textbf{\te{ఒ}} first, the way you learned it, then copy the small hook at its top.

\begin{quote}
This book does not yet say where the pen starts on this shape, or which way it travels. That is taught with real variation from school to school, and it waits here until it can be written down from a source. Copying the printed shape needs no such source, and it puts the shape in your hand in the meantime.
\end{quote}

\subsection*{Your turn}
\begin{itemize}
\raggedright
  \item \emph{Find:} \te{ఓ} at the start of \textbf{\te{ఓటు}}
  \item \emph{Trace it:} \te{ఓ} three times, saying \emph{ō} as you finish each one
  \item \emph{Write it:} \te{ఒ}, then \te{ఓ}, from memory, naming each
  \item \emph{Write it:} \textbf{\te{ఓటు}}
\end{itemize}

\begin{tcolorbox}[breakable,colback=teal!4,colframe=teal!35!black,title={Before you move on}]
Which letter is this — \te{ఓ}? (\textbf{ō}, long \emph{ō} at the start of a word.) Which word did it come from? (\textbf{\te{ఓటు}}, \emph{ōṭu}, \textquotedblleft{}a vote\textquotedblright{}.) How is it different from \textbf{\te{ఒ}}? (\textbf{The small hook at its top; the vowel is long.})
\end{tcolorbox}
\section[\texorpdfstring{\te{నడవడం} \te{వల్ల} (naḍavaḍaṁ valla)}{నడవడం వల్ల (naḍavaḍaṁ valla)}]{\te{నడవడం} \te{వల్ల} (naḍavaḍaṁ valla) — giving the reason}

\label{lesson:TE-C319-nadavadam-valla}



\begin{quote}
Say \textquotedblleft{}because of\textquotedblright{} (\emph{valla}). Then say \textquotedblleft{}a vote\textquotedblright{}. (\emph{ōṭu}.) Which letter opens it? (\textbf{\te{ఓ}}, a long \emph{ō}.)
\end{quote}

\subsection*{You'll want to know: -\te{అడం} \te{వల్ల}}
\textbf{-\te{అడం}} turns a verb into a noun, \textquotedblleft{}-ing\textquotedblright{}: \textbf{\te{నడవడం}}, \textquotedblleft{}walking\textquotedblright{}; \textbf{\te{చదవడం}}, \textquotedblleft{}reading\textquotedblright{}; \textbf{\te{తినడం}}, \textquotedblleft{}eating\textquotedblright{}; \textbf{\te{రావడం}}, \textquotedblleft{}coming\textquotedblright{}. With \textbf{\te{వల్ల}}, \textquotedblleft{}because of\textquotedblright{}, it gives a reason:

\begin{quote}
\textbf{\te{ప్రతి} \te{రోజు} \te{నడవడం} \te{వల్ల} \te{ఆరోగ్యం} \te{బాగుంటుంది}.}
\end{quote}

\begin{quote}
\emph{prati rōju naḍavaḍaṁ valla ārōgyaṁ bāguṇṭundi.} — \textquotedblleft{}Walking every day keeps you healthy.\textquotedblright{}
\end{quote}

\begin{quote}
\textbf{\te{ఎక్కువ} \te{తినడం} \te{వల్ల} \te{నాకు} \te{కడుపు} \te{నొప్పి} \te{వచ్చింది}.}
\end{quote}

\begin{quote}
\emph{ekkuva tinaḍaṁ valla nāku kaḍupu noppi vaccindi.} — \textquotedblleft{}Eating too much gave me a stomach ache.\textquotedblright{}
\end{quote}

\textbf{\te{బాగుంటుంది}} (\emph{bāguṇṭundi}, \textquotedblleft{}it stays well\textquotedblright{}) is \textbf{\te{బాగు}}, \textquotedblleft{}good, well\textquotedblright{}, the base of \textbf{\te{బాగా}}, with \textbf{\te{ఉంటుంది}}, \textquotedblleft{}it stays\textquotedblright{}, from \textbf{\te{ఉండు}}.

\begin{grammarlens}[title={Grammar Lens: two places for a reason}]
\begin{quote}
\textbf{\te{నేను} \te{రాలేదు}, \te{ఎందుకంటే} \te{నాకు} \te{జ్వరం} \te{వచ్చింది}.}
\end{quote}

\begin{quote}
\emph{nēnu rālēdu, endukaṇṭē nāku jvaraṁ vaccindi.} — \textquotedblleft{}I did not come, because I had a fever.\textquotedblright{}
\end{quote}

\begin{quote}
\textbf{\te{నాకు} \te{జ్వరం} \te{వచ్చింది} \te{కాబట్టి} \te{నేను} \te{రాలేదు}.}
\end{quote}

\begin{quote}
\emph{nāku jvaraṁ vaccindi kābaṭṭi nēnu rālēdu.} — \textquotedblleft{}I had a fever, so I did not come.\textquotedblright{}
\end{quote}

The same two facts. With \textbf{\te{ఎందుకంటే}} the reason stands after the result, where English puts it. With \textbf{\te{కాబట్టి}}, or \textbf{-\te{అడం} \te{వల్ల}}, it stands before, and the sentence ends on its main verb, as a Telugu sentence does. Both are everyday Telugu.
\end{grammarlens}

\subsection*{Your turn}
Say these aloud:

\begin{itemize}
\raggedright
  \item \textquotedblleft{}Walking every day keeps you healthy\textquotedblright{} — \emph{prati rōju naḍavaḍaṁ valla ārōgyaṁ bāguṇṭundi}
  \item \textquotedblleft{}Eating too much gave me a stomach ache\textquotedblright{} — \emph{ekkuva tinaḍaṁ valla nāku kaḍupu noppi vaccindi}
  \item \textquotedblleft{}I did not come, because I had a fever\textquotedblright{} — \emph{nēnu rālēdu, endukaṇṭē nāku jvaraṁ vaccindi}
  \item \textquotedblleft{}the same, with the reason first\textquotedblright{} — \emph{nāku jvaraṁ vaccindi kābaṭṭi nēnu rālēdu}
  \item why everyone must vote, in one sentence of your own
\end{itemize}

\begin{tcolorbox}[breakable,colback=teal!4,colframe=teal!35!black,title={Before you move on}]
\textquotedblleft{}Reading\textquotedblright{}, as a noun? (\textbf{\te{చదవడం}}.) \textquotedblleft{}Because of reading\textquotedblright{}? (\textbf{\te{చదవడం} \te{వల్ల}}.) With which word does the reason come after the result? (\textbf{\te{ఎందుకంటే}}.)
\end{tcolorbox}
\section[(review)]{Review — advice and reasons}

\label{lesson:TE-R319-advice-and-reasons}



\begin{quote}
Say \textquotedblleft{}a vote\textquotedblright{}. (\emph{ōṭu}.)

\begin{itemize}
\raggedright
  \item \emph{Write it:} \textbf{\te{ఓ}} once from memory; then \textbf{\te{ఈ}}, from \textbf{\te{ఈలోగా}}
\end{itemize}
\end{quote}

\subsection*{Your turn}
\begin{itemize}
\raggedright
  \item \emph{Say it:} You must not go to school today. You had better stay at home.
  \item \emph{Say it:} I think everyone must vote.
  \item \emph{Say it:} Because of reading every day, my Telugu is good.
  \item \emph{Say it:} I did not come, because I had a fever.
  \item \emph{Say it:} I had a fever. That is why I did not come.
  \item \emph{Write it:} \textbf{\te{ఓటు}}, then \textbf{\te{ఈ} \te{ఓటు}}, \textquotedblleft{}this vote\textquotedblright{}
\end{itemize}

\begin{tcolorbox}[breakable,colback=teal!4,colframe=teal!35!black,title={Before you move on}]
\textquotedblleft{}Must\textquotedblright{}? (\textbf{-\te{ఆలి}}.) \textquotedblleft{}Must not\textquotedblright{}? (\textbf{-\te{అకూడదు}}.) \textquotedblleft{}Had better\textquotedblright{}? (\textbf{-\te{తే} \te{మంచిది}}.) \textquotedblleft{}Because of walking\textquotedblright{}? (\textbf{\te{నడవడం} \te{వల్ల}}.) \textquotedblleft{}A vote\textquotedblright{}? (\textbf{\te{ఓటు}}, which opens with \textbf{\te{ఓ}}.)
\end{tcolorbox}
